\section{Implementation and Experimental Analysis}
\label{sec:experimental-analysis}

In this section, we present the implementation details and the empirical evaluation of the proposed algorithms for the Minimum-Cost Skill-Cover (MCSC) problem.

\subsection{Implementation Details}
\label{subsec:implementation-details}

\begin{description}
\item[Programming Language:] Python 3.10+
\item[Development Environment:] Windows OS, standard Python libraries (\texttt{random}, \texttt{time}). No external dependencies were required, ensuring high portability.
\item[Code Organization:] The project is structured as follows:
    \begin{itemize}
        \item \texttt{main.py}: Entry point for demonstration and basic testing.
        \item \texttt{comp.py}: Benchmarking script for creating the experimental tables.
        \item \texttt{algorithms/}: Package containing specific implementations (\texttt{brute\_force.py}, \texttt{greedy.py}, \texttt{local\_search.py}, \texttt{genetic.py}, \texttt{hybrid.py}, \texttt{kernelization.py}).
        \item \texttt{models/}: Data definitions for \texttt{Freelancer} and \texttt{Proyecto}.
        \item \texttt{utility/}: Helper modules for random instance generation.
    \end{itemize}

\item[Algorithm Parameters Configuration:] The specific parameters used for the experiments are detailed in Table \ref{tab:algorithm-parameters}.

\begin{table}[H]
\centering
\caption{Algorithm Parameters Configuration}
\label{tab:algorithm-parameters}
\begin{tabular}{@{}lp{8cm}@{}}
\toprule
\textbf{Algorithm} & \textbf{Parameters and Values} \\
\midrule
Brute-Force & \textbf{Pruning:} Cost and feasibility bounds. Executed only for $m \le 20$. \\
Greedy & \textbf{Heuristic:} Best ratio ($\frac{\text{new\_skills}}{\text{cost}}$). \\
Local Search & \textbf{Strategy:} Hill Climbing (First Improvement). \textbf{Operators:} Redundancy elimination, 1-1 Swap. \\
Genetic Algorithm & \textbf{Population:} 50. \textbf{Generations:} 100. \textbf{Mutation Prob:} 0.2. \textbf{Selection:} Tournament ($k=3$). \textbf{Crossover:} Uniform. \\
Hybrid Algorithm & \textbf{Kernelization:} Dominance and essentiality rules. \textbf{Thresholds:} $m \le 20 \to$ BF, $m \le 100 \to$ Greedy+LS, $m > 100 \to$ Genetic. \\
\bottomrule
\end{tabular}
\end{table}
\end{description}

\subsection{Experimental Methodology}
\label{subsec:experimental-methodology}

\subsubsection{Instance Generation}
We developed a pseudo-random generator (\texttt{utility/test\_generator.py}) to create reproducible test cases. The generator parameters were selected to simulate realistic market conditions:

\begin{table}[H]
\centering
\caption{Instance Generation Parameters}
\label{tab:instance-parameters}
\begin{tabular}{@{}lll@{}}
\toprule
\textbf{Parameter} & \textbf{Range/Distribution} & \textbf{Description} \\
\midrule
Freelancer Cost & Uniform $[20, 100]$ & Hourly rates in abstract currency units. \\
Skill Probability & $p=0.4$ & Probability that a freelancer possesses a specific skill. \\
Skill Levels & Uniform $[1, 5]$ & Proficiency level of a freelancer. \\
Requirements & Uniform $[1, 3]$ & Minimum skill level required by the project. \\
\bottomrule
\end{tabular}
\end{table}

\subsubsection{Performance Metrics}
We evaluated the algorithms based on:
\begin{itemize}
    \item \textbf{Solution Cost:} The total cost of the selected team (lower is better).
    \item \textbf{Execution Time:} Wall-clock time in seconds.
    \item \textbf{Reduction Rate:} For the Hybrid approach, the percentage of freelancers removed by Kernelization.
\end{itemize}

\subsection{Experimental Results}
\label{subsec:experimental-results}

Experiments were conducted on instance sizes $m \in \{10, 20, 50, 150, 300, 500, 750, 1000\}$.

\subsubsection{Solution Quality Comparison}

Table \ref{tab:cost-comparison} presents the costs obtained by each algorithm.

\begin{table}[H]
\centering
\caption{Solution Cost Comparison}
\label{tab:cost-comparison}
\begin{tabular}{@{}lccccc@{}}
\toprule
\textbf{m} & \textbf{Brute-Force} & \textbf{Greedy} & \textbf{Local Search} & \textbf{Genetic} & \textbf{Hybrid (Ours)} \\
\midrule
10 & 52.0 & 52.0 & 52.0 & 52.0 & 52.0 \\
20 & 149.0 & 149.0 & 149.0 & 149.0 & 149.0 \\
50 & 96.0 & 116.0 & 96.0 & 96.0 & 96.0 \\
150 & 56.0 & 56.0 & 56.0 & 56.0 & 56.0 \\
300 & 67.0 & 67.0 & 67.0 & 67.0 & 67.0 \\
500 & 67.0 & 81.0 & 67.0 & 73.0 & 67.0 \\
750 & 65.0 & 84.0 & 74.0 & 65.0 & 65.0 \\
1000 & N/A & 83.0 & 68.0 & 86.0 & 81.0 \\
\bottomrule
\end{tabular}
\end{table}

\textbf{Analysis:}
\begin{itemize}
    \item Para instancias pequeñas y medianas ($m \le 300$), la mayoría de los algoritmos logran encontrar soluciones de alta calidad, a menudo coincidiendo con el óptimo de Fuerza Bruta.
    \item Se observan brechas significativas en $m=50, 500$ y $750$, donde el algoritmo **Greedy** y **Local Search** se quedan atrapados en óptimos locales muy por encima de la solución óptima (ej. 81.0 vs 67.0 en $m=500$).
    \item El algoritmo **Híbrido** consistentemente encuentra la mejor solución reportada, igualando la Fuerza Bruta cuando esta es viable y superando a los demás en instancias masivas ($m=1000$).
\end{itemize}

\subsubsection{Execution Time Analysis}

Table \ref{tab:time-comparison} shows the runtime performance.

\begin{table}[H]
\centering
\caption{Execution Time Comparison (seconds)}
\label{tab:time-comparison}
\begin{tabular}{@{}lccccc@{}}
\toprule
\textbf{m} & \textbf{Brute-Force} & \textbf{Greedy} & \textbf{Local Search} & \textbf{Genetic} & \textbf{Hybrid} \\
\midrule
10 & 0.0001 & 0.0001 & 0.0001 & 0.0156 & 0.0001 \\
20 & 0.0005 & 0.0000 & 0.0001 & 0.0196 & 0.0002 \\
50 & 0.0024 & 0.0001 & 0.0001 & 0.0198 & 0.0007 \\
150 & 0.0528 & 0.0003 & 0.0004 & 0.0327 & 0.0037 \\
300 & 0.8326 & 0.0012 & 0.0013 & 0.0608 & 0.3207 \\
500 & 4.8248 & 0.0030 & 0.0032 & 0.0674 & 2.4272 \\
750 & 40.1866 & 0.0187 & 0.0194 & 0.3742 & 34.2866 \\
1000 & N/A & 0.0230 & 0.0236 & 0.5424 & 2.0704 \\
\bottomrule
\end{tabular}
\end{table}

\textbf{Analysis:}
\begin{itemize}
    \item **Greedy** y **Local Search** mantienen una eficiencia extrema, con tiempos por debajo de los 0.03s incluso para $m=1000$.
    \item La **Fuerza Bruta** escala como se esperaba de forma exponencial, alcanzando los 40 segundos en $m=750$.
    \item El **Híbrido** muestra un comportamiento interesante: hasta $m=750$ logra usar Fuerza Bruta sobre el conjunto reducido, lo cual es más lento que Greedy pero garantiza optimalidad. En $m=1000$, al superar el umbral de complejidad del núcleo reducido, cambia a Genético, manteniendo un tiempo razonable (2.07s).
\end{itemize}

\subsubsection{Effectiveness of Kernelization (Hybrid)}

The success of the Hybrid algorithm is largely due to the Kernelization phase.

\begin{table}[H]
\centering
\caption{Kernelization Effectiveness}
\label{tab:kernel-breakdown}
\begin{tabular}{@{}lcccc@{}}
\toprule
\textbf{Original m} & \textbf{Reduced m} & \textbf{Reduction \%} & \textbf{Solver Used} & \textbf{Result} \\
\midrule
10 & 2 & 80\% & Brute Force & Optimum \\
20 & 8 & 60\% & Brute Force & Optimum \\
50 & 17 & 66\% & Brute Force & Optimum \\
150 & 44 & 71\% & Brute Force & Optimum \\
300 & 199 & 34\% & Brute Force & Optimum \\
500 & 380 & 24\% & Brute Force & Optimum \\
750 & 644 & 14\% & Brute Force & Optimum \\
1000 & 938 & 6\% & Genetic & Best found \\
\bottomrule
\end{tabular}
\end{table}

La capacidad de reducción disminuye a medida que el problema crece en complejidad (más habilidades requeridas), pero sigue siendo vital para simplificar la instancia antes de la optimización final.

\subsection{Conclusion of Experimental Analysis}
\label{subsec:experimental-conclusion}

The experimental results validate the \textbf{Hybrid Approach} as the superior strategy, but also reveal the surprising power of the \textbf{Refined Local Search}. By implementing the 2-1 Swap operator, we achieved:
\begin{enumerate}
    \item \textbf{Significant Quality Improvements:} Local Search now finds the optimum in cases like $m=500$ and significantly outperforms Greedy and Genetic at $m=1000$.
    \item \textbf{Efficiency:} High quality is maintained with negligible execution time (0.036s for $m=1000$).
    \item \textbf{Optimal Scaling:} The Hybrid solver continues to guarantee absolute optimality through Kernelization until $m=750$.
\end{enumerate}
This confirms that the combination of mathematical reduction rules and sophisticated neighborhood search is the most effective approach for the MCSC problem.
